\begin{remark}{2.12}\label{XIV.2.12}
I do not know whether statement 2.11 a) and b) is valid without a
hypothesis of existence of regular points (when $k(s)$ is finite).
\end{remark}
It remains to prove implication $(\mathrm C_3)\Rightarrow(\mathrm C_2)$.
Note also the following equivalent form of $(\mathrm C_3)$:

\noindent $(\mathrm C'_3)$ One has $(\mathrm C_0)$, i.e. the nilpotent
rank of the $\mathfrak g(s)$ ($s\in S$) is locally constant, and on
every open subset $V$ of $S$ where this rank has value $r$, the Killing
polynomial of $\mathfrak g|_V$ is divisible by $t^r$.

One must show that this condition implies that every quasi-regular
section $a$ of $\mathfrak g$ is regular. Taking 2.6 into account,
this is contained in the following lemma (applied to the endomorphism
$\ad(a)$ of $\mathfrak g$), (iv) $\Rightarrow$ (iii):

\begin{lemma}{2.13}\label{XIV.2.13}
Let $A$ be a ring, $M$ a projective $A$-module of finite type, and
$u$ an endomorphism of $M$. The following conditions are equivalent:
\begin{enumeratei}
\item $M$ is the direct sum of two stable modules $M',M''$ such that
$u|_{M'}$ is nilpotent, and $u|_{M''}$ is an automorphism of $M''$.
\item There exists an integer $n>0$ such that $\Im u^n+\Ker u^n=M$.
\item The nilspace $\operatorname{Nil}(u)=\bigcup_{n>0}\Ker u^n$ is a
direct summand of $M$, and $M=\operatorname{Nil}(u)+u(M)$.
\end{enumeratei}
These conditions are implied by the following one (and are equivalent
to it when $A$ is reduced):

\noindent (iv) Locally on $\Spec(A)$ (for the Zariski topology) the
characteristic polynomial $P_u(t)$ of $u$ can be put in the form
\[
 t^r(t^{n-r}+c_1t^{n-r-1}+\cdots+c_{n-r}),
\]
where $c_{n-r}$ is invertible.
\end{lemma}
The equivalence of (i), (ii), (iii) is immediate and is recorded for
reference. The fact that (i) implies (iv) when $A$ is reduced follows
from the fact that in this case a nilpotent endomorphism of a projective
module of rank $r$ has characteristic polynomial $t^r$, whereas in
all cases the characteristic polynomial of an automorphism of a
projective module of finite type has as constant term the determinant
of $u$ up to sign (locally on $\Spec(A)$), hence an invertible element
of $A$. Finally, to prove (iv) $\Rightarrow$ (i), one notes that $M$
is a module over the polynomial ring $A[t]$, letting $t$ act by $u$,
and the well-known identity
\[
 P(u)=0
\]
shows that $M$ is annihilated by $PA[t]$, hence can be considered as
a module over $A[t]/PA[t]$. Now, writing $P=t^rQ$, where the constant
term of $Q$ is invertible, one sees at once that
\[
 PA[t]=t^rA[t]\cap QA[t],
\]
so $A[t]/PA[t]$ decomposes as the product of the rings $A[t]/t^rA[t]$
and $A[t]/QA[t]$, whence a corresponding decomposition of $M$ as
a sum of two $A[t]$-modules, i.e. as a sum of two sub-$A$-modules
$M'$ and $M''$ stable under $u$; this is the decomposition considered
in (i).

This completes the proof of 2.13, hence of 2.9.
\end{proof}

\begin{corollary}{2.14}\label{XIV.2.14}
When condition $(\mathrm C_3)$ of 2.9 holds, the Cartan subalgebras
of $\mathfrak g$ are strictly nilpotent (2.2).
\end{corollary}
The proof is immediate.

\begin{remark}{2.15}\label{XIV.2.15}
\noindent a) Let us point out that one can prove a converse of 2.14:
$(\mathrm C_3)$ is equivalent to the fact that for every $S'$ over
$S$, every quasi-regular section of $\mathfrak g_{S'}$ is regular
and every Cartan subalgebra of $\mathfrak g_{S'}$ strictly nilpotent,
or again every quasi-regular section of $\mathfrak g_{S'}$ contained
in a strictly nilpotent Cartan subalgebra of $\mathfrak g_{S'}$.

\noindent b) Unlike the other conditions $(\mathrm C_0)$ to
$(\mathrm C_2)$, condition $(\mathrm C_3)$ is of ``infinitesimal
nature''; more precisely, when $S$ is locally noetherian,
$\mathfrak g$ satisfies condition $(\mathrm C_3)$ if and only if
for every base change $S'\to S$, with $S'$ artinian local (if one
wishes, $S'$ the spectrum of an artinian quotient of a local ring
of $S$), $\mathfrak g_{S'}$ satisfies the same condition. Likewise,
when $(\mathrm C_0)$ holds, the condition for a section of
$\mathfrak g$ to be regular is of infinitesimal nature.

\noindent c) When $\mathfrak g$ is the Lie algebra of a smooth group
prescheme of finite presentation over $S$, then we shall see that
conditions $(\mathrm C_0),(\mathrm C_1),(\mathrm C'_1),(\mathrm C_2)$
on $\mathfrak g$ are equivalent (5.2 a)); I do not know what happens
in general (except that, even for $S$ artinian local,
$(\mathrm C_0)$ does not imply $(\mathrm C_1)$). However, even in
the case where $\mathfrak g$ comes from a $G$, and $S$ is artinian
local, it is not true in general that $(\mathrm C_2)$ implies
$(\mathrm C_3)$, for $\mathfrak g$ can be nilpotent without being
strictly nilpotent. One obtains an example of this fact by starting
from a smooth affine group scheme $G$ over the spectrum $S$ of a
discrete valuation ring, such that the Lie algebra of the general
fiber is not nilpotent (for example, the general fiber is an adjoint
semisimple group), and that of the special fiber is nilpotent
(for example, the special fiber being a vector group): then for
$n$ sufficiently large, the Lie algebra of $G_n=G\times_S S_n$
is not strictly nilpotent, but it is nilpotent.
\end{remark}

Let us still place ourselves under the conditions of 2.4, and let
\[
 \mathscr D:(\Sch)^\circ_{/S}\longrightarrow\Ens
\]
be the functor defined by
\[
 \mathscr D(S')=\text{set of Cartan subalgebras of }\mathfrak g_{S'}.
\]
Introduce also the functor
\[
\begin{aligned}
 X(S')=\{&(\mathfrak d,a):\mathfrak d\text{ is a Cartan subalgebra of }\mathfrak g_{S'}\\
 &\text{and }a\text{ a section of }\mathfrak d\}.
\end{aligned}
\]
One thus has two projections $(\mathfrak d,a)\mto\mathfrak d$ and
$(\mathfrak d,a)\mto a$:
\[
 p:X\longrightarrow\mathfrak d\qquad\text{and}\qquad
 \psi:X\longrightarrow\mathrm W(\mathfrak g).
\]
When $(\mathrm C_0)$ holds, we also consider the open subset $U$
of regular points of $\mathrm W(\mathfrak g)$ (cf. 2.10), and
condition $(\mathrm C_2)$ is then expressed by the fact that the
morphism
\[
 \psi^{-1}(U)\longrightarrow U
\]
induced by $\psi$ is an isomorphism (a priori, it is a monomorphism
by 2.6). Note that it is trivial that the morphism
$p:X\to\mathfrak d$ is representable by a projection of vector
bundles (i.e. for every $S$-morphism $S'\to\mathfrak d$,
corresponding to a Cartan subalgebra $\mathfrak d$ of
$\mathfrak g_{S'}$, $X\times_{\mathfrak d}S'$ is representable
by a vector bundle over $S'$, namely $\mathrm W(\mathfrak d)$);
hence if $\mathfrak d$ is representable, so is $X$. Now one has:
% typo?: kept as printed: fraktur d is used as the target instead of script D.

\begin{theorem}{2.16}\label{XIV.2.16}
Let $S$ be a prescheme, $\mathfrak g$ a Lie algebra over $S$
which is a locally free $\OS$-module of finite type; suppose
condition $(\mathrm C_0)$ of 2.9 holds.

\noindent a) The functor $\mathscr D$ of Cartan subalgebras of
$\mathfrak g$ defined above is representable by a quasi-projective
prescheme of finite presentation over $S$. The same holds for
the functor $X$ defined above.

\noindent b) When condition $(\mathrm C_2)$ of 2.9 holds,
$\mathscr D$ and $X$ are smooth over $S$, and the morphism
$\psi^{-1}(U)\to U$ induced by $\psi$ is an isomorphism.

\noindent c) Still supposing condition $(\mathrm C_2)$ holds,
let $s\in S$, and let $\mathfrak d_0$ be a Cartan subalgebra
of $\mathfrak g(s)$ corresponding to a point $d$ of
$\mathscr D$ rational over $k(s)$. Suppose that
$\mathfrak d_0$ contains a regular point of $\mathfrak g(s)$
(a condition automatically verified if $k(s)$ is infinite).
Let $r$ be the infinitesimal rank of $\mathfrak g(s)$,
and $n$ its rank over $k(s)$; then there exists an open
neighborhood $V$ of $d$ in $\mathscr D$ which is
$S$-isomorphic to an open subset $V'$ of $S[t_1,\ldots,t_{n-r}]$.
\end{theorem}
\begin{proof}
One may suppose that $\mathfrak g$ has constant rank $n$
and constant nilpotent rank $r$. The assertions made about
$X$ in a) and b) follow at once from the assertions made
about $\mathscr D$ and the fact that $X$ is a vector bundle
over $\mathscr D$ defined by a locally free module, and
are included merely for reference.

\noindent a) The functor $\mathscr D$ is a subfunctor of
the functor $\underline{\operatorname{Grass}}_{n-r}(\mathfrak g)$
whose value at $S'$ is the set of locally free quotient
modules of rank $n-r$ of $\mathfrak g_{S'}$, and it is
well known that the latter functor is representable by
a projective prescheme smooth over $S$ (cf., for example,
S\'eminaire Cartan 1960/61, Exp. 12, Nos. 2 and 3, whose
constructions carry over unchanged to the case of
preschemes)%
{\renewcommand{\thefootnote}{*}\footnote{Cf. also EGA I,
2nd edition (to appear in North Holland Publishing Cie).}\addtocounter{footnote}{-1}}%
\nde{See \S I.1.3 of M. Demazure and P. Gabriel,
\emph{Groupes alg\'ebriques}, Masson (1970).}.
One is therefore reduced to a relative problem, namely
the following: given a locally free quotient module
of rank $n-r$ of $\mathfrak g$, or, what amounts to
the same thing, a locally free submodule $\mathfrak d$
of rank $r$ which is locally a direct summand, represent
the following functor: $F(S')=\varnothing$ if
$\mathfrak d$ is not a Cartan subalgebra of
$\mathfrak g_{S'}$, $F(S')=\{\varnothing\}$ otherwise.
In fact, we shall see that $F$ is representable by
a subprescheme of finite presentation of $S$ (which
will show that $\mathscr D\to\underline{\operatorname{Grass}}$
is representable by an immersion of finite presentation,
and will finish proving a)). One begins by expressing
the condition that $\mathfrak d_{S'}$ be a Lie
subalgebra of $\mathfrak g_{S'}$; one sees immediately
that this is expressed by the fact that $S'\to S$
factors through a certain closed subprescheme $S_1$
of $S$, of finite presentation over $S$ (whose local
equations over $S$ are written down immediately using
a basis of $\mathfrak g$ adapted to the submodule
$\mathfrak d$). One may therefore suppose that one
already has $S=S_1$. One must next express that
$\mathfrak d_{S'}$ contains, locally for fpqc, a
quasi-regular section of $\mathfrak g_{S'}$, and for
this one considers $V=\mathrm W(\mathfrak d)\cap U$,
where $U$ is the open subset of regular points of
$\mathrm W(\mathfrak g)$ (2.10); then the structural
morphism $V\to S$ being smooth and quasi-compact,
its image $S_1$ is an open part of $S$, and the
immersion morphism $S_1\to S$ is quasi-compact,
i.e. of finite presentation. The condition considered
on $S'$ is then expressed by saying that $S'\to S$
factors through $S_1$. Thus one is reduced to the
case $S=S_1$, and, using descent theory, to the
case where $\mathfrak d$ admits a section $a$
which is a quasi-regular section of $\mathfrak g$.
One must finally express that the section $a_{S'}$
of $\mathfrak g_{S'}$ deduced from $a$ satisfies
$\ad(a_{S'})_{\mathfrak g_{S'}/\mathfrak d_{S'}}$
bijective, which again amounts to saying that
$S'\to S$ factors through a certain open subprescheme
of finite presentation of $S$, namely $S_D$, where
$D$ is the determinant of $\ad(a)_{\mathfrak g/\mathfrak d}$.
But then one sees immediately that $\mathfrak d|_{S_D}$
is a Cartan subalgebra of $\mathfrak g|_{S_D}$,
so $S_D$ represents the functor $F$, which proves a).

\noindent b) Is immediate by 2.11 a) and XI 1.5.
Of course, b) is also a consequence of the more
precise statement c).

\noindent c) Let $a_0$ be a regular point of
$\mathfrak g(s)$ contained in $\mathfrak d_0$;
extend it to a section $a$ of $\mathfrak g$
over an open neighborhood $V$ of $s$; one may
clearly suppose $V=S$. On the other hand, let
$M_0$ be a complement of the vector space
$\mathfrak d_0$ in $\mathfrak g(s)$; then in
an open neighborhood $V$ of $S$ there exists
a submodule $M$ of $\mathfrak g$, a direct
summand of $\mathfrak g|_V$, such that $M(s)=M_0$,
and one may again suppose $V=S$. Let now $\mathbf V$
be the subfunctor of $\mathscr D$ such that
$\mathbf V(S')$ is the set of Cartan subalgebras
$\mathfrak d'$ of $\mathfrak g_{S'}$ satisfying
the following two conditions:
\begin{enumerate}[label=$\arabic*^\circ)$]
\item $\mathfrak d'$ is complementary to $M_{S'}$, and
\item the unique section of $(a_{S'}+M_{S'})\cap\mathfrak d$
is a regular section of $\mathfrak g_{S'}$.
\end{enumerate}
Condition $1^\circ)$ corresponds to an open subset
$V_1$ of $\mathfrak d$ (induced by the open subset
of $\underline{\operatorname{Grass}}_{n-r}(\mathfrak g)$
defined by the same condition $1^\circ)$); the
conjunction of $1^\circ$ and $2^\circ$ corresponds
to an open subset of $V_1$ by 2.10 and
$(\mathrm C_2)$. Thus $\mathbf V$ is represented
by an open subprescheme $V$ of $\mathscr D$,
clearly containing $a$.
% typo?: kept as printed: d rather than d' in condition 2; V contains a.

On the other hand, let $\mathbf V'$ be the
subfunctor of $\mathrm W(M)$ defined by
\[
\begin{aligned}
 \mathbf V'(S')=\{&\text{sections }u'\text{ of }M_{S'}\text{ such that:}\\
 &\text{(i) }a_{S'}+u'\text{ is a regular section of }\mathfrak g_{S'},\text{ and}\\
 &\text{(ii) the unique Cartan subalgebra }\mathfrak d'\text{ of }\mathfrak g_{S'}\\
 &\text{containing }a_{S'}+u'\text{ is complementary to }M_{S'}\}.
\end{aligned}
\]
Then condition $1^\circ)$ corresponds to an open
subprescheme $V'_1$ of $\mathrm W(M)$, namely
the inverse image of the open subset $U$ of
regular points of $\mathrm W(\mathfrak g)$
(cf. 2.10) under the translation morphism
$m\mto a+m$. The conjunction of (i) and (ii)
corresponds to an open subset $V'$ of $V'_1$,
namely the inverse image of $V$ under the
obvious morphism $V'_1\to\mathscr D$
(associating to $u'$ the unique Cartan
subalgebra $\mathfrak d'$ of $\mathfrak g_{S'}$
containing $a_{S'}+u'$). The restriction of
this last morphism to $V'$ is a morphism
\[
 V'\longrightarrow V,
\]
which is clearly an isomorphism. This proves c).
% typo?: kept as printed: condition 1 degrees rather than (i).
\end{proof}

\begin{corollary}{2.17}\label{XIV.2.17}
Let $\mathfrak g$ be a finite-dimensional Lie
algebra over a field $k$. Then the scheme
$\mathscr D$ of Cartan subalgebras of
$\mathfrak g$ (2.16 a)) is quasi-projective,
smooth and irreducible. When $\mathfrak g$
contains a regular element (for example
when $k$ is infinite), $\mathfrak d$ is a
rational variety, i.e. its function field
is a pure extension of $k$.
\end{corollary}
The fact that $\mathscr D$ is irreducible
follows from the fact that one has a
surjective morphism $\psi^{-1}(U)\to\mathscr D$,
and $\psi^{-1}(U)$ is irreducible, being
isomorphic to the open subset $U$ of
$\mathrm W(\mathfrak g)$. The assertion
about the function field is an immediate
consequence of c).

\begin{remarks}{2.18}\label{XIV.2.18}
I do not know whether this conclusion
remains true if $k$ is finite, without
supposing that $\mathfrak g$ contains a
regular point; compare 2.12. One can prove
that it is so when $\mathfrak g$ is the
Lie algebra of a smooth algebraic group
$G$ over $k$, at least when $G/\text{radical}$
is an ``adjoint'' semisimple group, using
a result of Chevalley mentioned below
(cf. Appendix). It is plausible that this
result remains valid without restriction
on $G$; for this it would suffice that
the cited result of Chevalley be proved
for every semisimple algebraic group
(not necessarily adjoint).
\end{remarks}

\sgasection{3}{Subgroups of type (C) of group preschemes over an arbitrary prescheme}
\begin{theorem}{3.1}\label{XIV.3.1}
Let $S$ be a prescheme, $G$ a smooth
$S$-group prescheme, $\mathfrak g$ its
Lie algebra (which is a locally free
module of finite type over $S$), and
$\mathfrak h$ a Lie subalgebra of
$\mathfrak g$ which is (as a module)
locally a direct summand in
$\mathfrak g$, and such that for every
$s\in S$, the geometric fiber
$\mathfrak h_{\bar s}$ contains a
Cartan subalgebra of $\mathfrak g_{\bar s}$.
Let $a$ be a quasi-regular section of
$\mathfrak g$ (2.5). Then
\[
 M_a=\uTransp_G(a,\mathfrak h)
\]
(the subfunctor of $G$ whose points
with values in $S'$ are the
$g\in G(S')$ such that
$\ad(g)\cdot a_{S'}\in\Gamma(s',\mathfrak h_S)$)
is representable by a closed subprescheme\nde{One
will denote it by $\Transp_G(a,\mathfrak h)$.}
of $G$ smooth over $S$, whose structural
morphism to $S$ is surjective.
% typo?: kept as printed: Gamma(s',h_S).
\end{theorem}
Consider the canonical morphism
\[
 \varphi:G\underset S\times\mathrm W(\mathfrak h)
 \longrightarrow\mathrm W(\mathfrak g)
\]
given by $(g,x)\mto\ad(g)\cdot x$;
then $M_a$ is $S$-isomorphic to
$\varphi^{-1}(a)$, the inverse image
of $a$ (considered as a section of
$\mathrm W(\mathfrak g)$ over $S$)
under $\varphi$. It will therefore
suffice for the smoothness of $M_a$
that we show that $\varphi$ is smooth
at the points of $G\times_S\mathrm W(\mathfrak h)$
over $\Im(a)$; more generally,
$\varphi$ is smooth at every point
over a point of $\mathrm W(\mathfrak g)$
which is regular in its fiber
$\mathrm W(\mathfrak g(s))$ over $S$.
To see this, since the source and
target of $\varphi$ are smooth,
hence flat and locally of finite
presentation over $S$, one is
reduced to verifying it fiber by
fiber, which reduces one to the
case where $S$ is the spectrum of
an algebraically closed field,
$G$ thus being a locally algebraic
group over $k$, $\mathfrak h$ a
subalgebra of its Lie algebra
$\mathfrak g$ containing a Cartan
subalgebra of $\mathfrak g$, and
$a$ a regular point of
$\mathfrak g$. One may clearly
suppose (taking into account that
$\varphi$ is a $G$-morphism) that
the point considered in
$G\times\mathrm W(\mathfrak h)$
is of the form $(e,a)$. One may
clearly suppose $G$ connected,
hence of finite type over $k$,
but then our assertion is
precisely XIII 5.4. Moreover,
the fact that $M_a$ is a closed
subprescheme of $G$ (of finite
presentation over $S$) is trivial,
since $M_a$ is the inverse
image of $\mathrm W(\mathfrak h)$
under the morphism
$g\mto\ad(g)\cdot a$ from $G$
to $\mathrm W(\mathfrak g)$.
The surjectivity of the
structural morphism $M_a\to S$
also reduces to the case of
an algebraically closed base
field, but then $\mathfrak h$
contains a Cartan subalgebra
$\mathfrak d$ by hypothesis,
which is therefore conjugate
to the Cartan subalgebra
$\mathfrak d=\operatorname{Nil}(a,\mathfrak g)$
by the conjugacy theorem XIII
6.1 a), hence $\mathfrak h$
contains a conjugate of $a$.
This completes the proof of 3.1.
% typo?: kept as printed: both Cartan subalgebras are named d.

\begin{corollary}{3.2}\label{XIV.3.2}
Let $G,\mathfrak g$ be as in
3.1, with $G$ of finite type
over $S$, and let
$\mathfrak k,\mathfrak h$ be
two Lie subalgebras of
$\mathfrak g$, locally direct
summands (as modules); suppose
that one is under one of
the following two hypotheses:

\noindent a) For every $s\in S$,
the geometric fiber
$\mathfrak k_{\bar s}$ is
nilpotent and contains a
regular element of
$\mathfrak g_{\bar s}$;
the geometric fiber
$\mathfrak h_{\bar s}$
contains a Cartan subalgebra
of $\mathfrak g_{\bar s}$.

\noindent b) $\mathfrak k$
is a Cartan subalgebra of
$\mathfrak g$.

Then $\uTransp_G(\mathfrak k,\mathfrak h)$
is a closed subprescheme of
$G$ smooth over $S$; moreover,
in case a), its structural
morphism to $S$ is surjective.
% typo: sous-akgebres -> sous-algebres.
\end{corollary}
The fact that the transporter
is representable by a closed
subprescheme of finite
presentation of $G$ is
immediate, and is left to
the reader. Let us first
prove smoothness in case a).
Suppose first that there
exists a section $a$ of
$\mathfrak k$ which is
quasi-regular in
$\mathfrak g$. Then it
suffices to apply 3.1 and the

\begin{lemma}{3.3}\label{XIV.3.3}
Under the conditions of
3.2 a), if $a$ is a section
of $\mathfrak k$ quasi-regular
in $\mathfrak g$, then one has
\[
 \uTransp_G(\mathfrak k,\mathfrak h)
 =\uTransp_G(a,\mathfrak h).
\]
\end{lemma}
Indeed, taking the definitions
into account, this amounts
to showing that if $a$ is
moreover a section of
$\mathfrak h$, then one
has $\mathfrak k\subset\mathfrak h$.
Now, since by hypothesis
$\mathfrak k$ is locally
nilpotent, it follows from
2.1 that $\mathfrak k\subset
\operatorname{Nil}(a,\mathfrak g)$;
on the other hand,
$\operatorname{Nil}(a,\mathfrak g)\subset\mathfrak h$
for $\ad(a)_{\mathfrak g/\mathfrak h}$
is injective (being so
fiber by fiber by XIII
4.8 b)). Whence the conclusion.
In the general case,
one reduces to the case
where $S$ is affine
noetherian by the usual
standard procedure, then
to the case where $S$ is
artinian local (smoothness
being a property of
infinitesimal nature),
and by flat descent to
the case where its residue
field is infinite, hence
the fiber $\mathscr K_0$
admits an element regular
in $\mathfrak g_0$.
One lifts this element
to an element of
$\mathscr K=\Gamma(\mathfrak k)$,
which reduces one to
the preceding case.
Thus one has proved,
in case a), smoothness
of the transporter;
as for the fact that
its structural morphism
is surjective, it reduces
to the case where $S$ is
the spectrum of an
algebraically closed
field, hence where
$\mathscr K$ contains
a regular point of
$\mathfrak g$, and
one applies 3.3 and 3.1.

To prove b), one is
reduced by the definition
of smoothness (XI 1.1)
to proving that if $S$
is affine, $S_0$ a
subscheme defined by a
nilpotent quasi-coherent
Ideal $J$, and $g_0$ an
element of $G(S_0)$
transporting $\mathfrak k_0$
into $\mathfrak g_0$,
then $g_0$ lifts to
an element $g$ of
$G(S)$ transporting
$\mathfrak k$ into
$\mathfrak g$. Now
the hypothesis made
on $g_0$ implies that
one is under the
conditions of a),
already treated.
This completes the proof.
% typo?: kept as printed: transports k_0 into g_0, and k into g.

Of course, when in
3.2 b) $\mathfrak h$
satisfies the stronger
hypothesis of 3.1,
then (and only then)
the structural morphism
$\uTransp_G(\mathfrak d,\mathfrak h)\to S$
is surjective. Using
Hensel's lemma XI 1.10,
one concludes from
3.1 and 3.2:
% typo?: kept as printed: d rather than k.

\begin{corollary}{3.4}\label{XIV.3.4}
Under the conditions
of 3.1 for $G$ and
$\mathfrak h$, and
supposing $G$ of
finite type over $S$:
